\section{Limitations and Conclusion}

In this work, we examine a simple strategy of leveraging an LLM's self-consistency to train it via an RL framework. Our experiments on synthetic reasoning tasks with controllable difficulty levels demonstrate the promise of such a framework: not only can LLMs improve their performance on these tasks, but they can also improve the quality of their self-supervision for the next training step. Our analysis, however, also reveals the limitation of leveraging self-consistency as training reward: prolonged training under such a framework can lead to reward hacking and complete model collapse. Future investigations could explore how to develop more robust forms of verification, extend self-verification to other domains like coding, and curriculum learning strategies that expose the model to problems of only the right difficulty. Additionally, leveraging LLM-as-judges or generative verifiers to improve the training signal, or employing additional consistency regularization between the chain-of-thought and the final answer to mitigate reward hacking can be promising future research. We leave these and other promising directions for sustained self-improvement for future work.

\section*{Acknowledgements}
This work has greatly benefited from the use of Delta's advanced computing and data resource supported by the National Science Foundation (OAC 2005572) and the State of Illinois, as part of ACCESS-approved compute grants~\citep{access_compute}. We also appreciate the computing resources of Bridges-2~\citep{psc_computing} at Pittsburgh Supercomputing Center through ACCESS allocation CIS240901 from the Advanced Cyberinfrastructure Coordination Ecosystem: Services \& Support (ACCESS) program, which is supported by National Science Foundation grants \#2138259, \#2138286, \#2138307, \#2137603, and \#2138296. Overall, this project used ACCESS grants CIS240901, CIS230366, CIS250216, and CIS250428 for its compute resources. Moreover, FT was partially supported by the U.S. Army Futures Command under Contract No. W519TC-23-C-0030 during the project. Andrea Zanette was partially supported by the National Science Foundation under Grants CCF-2106778 and DMS-2134080.

The authors thank Brandon Pusateri, Jillian Lehosky, and Greg Bauer from ACCESS Support Staff for their incredible help at approving supplements and renewals for ACCESS compute grants throughout this project. Moreover, the work would not have finished so quickly without the help of Brett Bode from NCSA Delta Support Staff, who provided the authors with critical help in properly utilizing the Delta cluster.  FT gratefully acknowledges Daman Arora, Yuda Song, Yiding Jiang, Ruiqi Zhang, Zhaoyi Zhou, Guanning Zeng, Yutong He and other members of the Zanette, Russ, Auton, and AIRe lab for feedback and suggestions received on earlier versions of this work.
